\documentclass[12pt]{article}
\usepackage[T1]{fontenc}
\usepackage[utf8]{inputenc}
\usepackage{lmodern}
\usepackage{textcomp}
\usepackage{enumitem}
\usepackage{geometry}
\geometry{margin=1in}
\begin{document}
\begin{center}
Answer Key - Sociology - Undergraduate Level Assessment 13 \
\small (Answers are ordered exactly as in the question paper.)
\end{center}

\section*{Multiple Choice}
\begin{enumerate}[label=\arabic*.]
\item \textbf{Answer:} C
\\ \textit{Options:} A) criminals were socialized into an 'underworld' of crime; B) no act is intrinsically deviant; C) biological failings drove some people into crime; D) women were less likely to be arrested than men
\\ \textit{Note:} Correct option: C - biological failings drove some people into crime
\item \textbf{Answer:} B
\\ \textit{Options:} A) contrasted homosexuality with heterosexuality; B) distinguished between male and females as separate sexes; C) represented women's genitalia as underdeveloped versions of men's; D) argued for male superiority over women
\\ \textit{Note:} Correct option: B - distinguished between male and females as separate sexes
\item \textbf{Answer:} D
\\ \textit{Options:} A) a means of transporting money between different areas of a country; B) an empire with a bureaucratic administration but no political centre; C) an awareness of risks and dangers that affect the environment as a whole; D) a unit with a division of labour that extends across ethnic and cultural groups
\\ \textit{Note:} Correct option: D - a unit with a division of labour that extends across ethnic and cultural groups
\item \textbf{Answer:} C
\\ \textit{Options:} A) the state's power can be exercised through several different administrative structures; B) the ruling elite is composed of people from various class backgrounds; C) political parties must compete for the votes of 'consumers' in the electorate; D) there is a close alignment between class background and party preference
\\ \textit{Note:} Correct option: C - political parties must compete for the votes of 'consumers' in the electorate
\item \textbf{Answer:} D
\\ \textit{Options:} A) domesticating its content, including more 'home-produced' programmes; B) controlling the distribution of imported products by banning satellite dishes; C) creating 'reverse flows' of their own programmes back to imperial societies; D) all of the above
\\ \textit{Note:} Correct option: D - all of the above
\end{enumerate}

\section*{True / False}
\begin{enumerate}[label=\arabic*.]
\setcounter{enumi}{5}
\item \textbf{Answer:} False
\\ \textit{Options:} A) True; B) False
\\ \textit{Note:} LLM-generated false statement. Correct answer: 'non-probability sampling'.
\item \textbf{Answer:} True
\\ \textit{Options:} A) True; B) False
\\ \textit{Note:} LLM-generated true statement based on MCQ.
\item \textbf{Answer:} True
\\ \textit{Options:} A) True; B) False
\\ \textit{Note:} LLM-generated true statement based on MCQ.
\item \textbf{Answer:} False
\\ \textit{Options:} A) True; B) False
\\ \textit{Note:} LLM-generated false statement. Correct answer: 'Auguste Comte'.
\item \textbf{Answer:} False
\\ \textit{Options:} A) True; B) False
\\ \textit{Note:} LLM-generated false statement. Correct answer: 'the media manipulate 'the masses' as vulnerable, passive consumers'.
\end{enumerate}

\section*{Long Form Response}
\begin{enumerate}[label=\arabic*.]
\setcounter{enumi}{10}
\item \textbf{Answer:} triangulation. Explain why this is the correct answer and provide supporting evidence. Reference key facts or concepts that support this choice.
\\ \textit{Note:} Long-form derived from MCQ; award credit for stating the correct idea and giving supporting reasoning.
\item \textbf{Answer:} freedom of movement for citizens. Explain why this is the correct answer and provide supporting evidence. Reference key facts or concepts that support this choice.
\\ \textit{Note:} Long-form derived from MCQ; award credit for stating the correct idea and giving supporting reasoning.
\end{enumerate}

\vfill
\noindent\textit{End of Answer Key}
\end{document}